\chapter{Introduction}
\label{chapter:introduction}

Low-thrust electric propulsion is attractive for deep-space missions because of its high specific
impulse, but it forces the trajectory design problem into a genuinely different regime than
chemical propulsion: instead of a handful of discrete impulsive burns, the spacecraft must be
steered continuously over months or years. The classical approach -- optimal control via
Pontryagin's minimum principle, or direct transcription and nonlinear programming -- produces
excellent open-loop solutions but is expensive to (re-)solve and is not naturally a feedback
policy. This raises a natural question for this project: can a policy trained with model-free
reinforcement learning (RL) instead learn a reusable, closed-loop \emph{guidance law} for
low-thrust rendezvous, and if it does, does the resulting control structure resemble the
bang-bang-like behaviour known to be optimal for mass/time-optimal low-thrust transfers?

\section{Problem statement}

The testbed is the rendezvous leg of the 4th Global Trajectory Optimization Competition (GTOC4):
a spacecraft with wet mass \SI{1500}{kg} (of which \SI{1000}{kg} is propellant), a constant
maximum thrust of \SI{0.135}{N} and specific impulse \SI{3000}{s}, must depart Earth and rendezvous
with one of 1436 catalogued asteroids, inside a launch window between 2015 and 2025 and a total
mission duration of up to 10 years. The problem is governed by two-body heliocentric dynamics
(Sun point-mass gravity plus a controllable low-thrust acceleration); the spacecraft's own
gravitational influence and other perturbations are neglected, matching the original GTOC4
problem statement.

\section{Scientific question}

This project asks two related questions:
\begin{enumerate}
  \item Can a policy-gradient RL agent (Proximal Policy Optimization, PPO \parencite{schulman2017ppo}) learn a low-thrust
        guidance law that rendezvous with a target under this dynamics model, starting from a
        curriculum of increasingly difficult targets?
  \item If it can, does the learned thrust profile recover the bang-bang-like structure expected
        from optimal control theory, and does the learned policy implement a genuine closed-loop
        feedback law, or does it merely memorise a single open-loop-like transfer?
\end{enumerate}

\section{Approach and tools}

The dynamics core is a hand-written, fixed-step 4th-order Runge-Kutta (RK4) integrator in NumPy,
chosen for speed inside the RL training loop; it is validated against Tudat's adaptive RKF78
integrator (used only offline, never inside the training loop). The environment is implemented
with the Gymnasium API and trained with Stable-Baselines3's PPO implementation (PyTorch, CPU-only).
A curriculum of three stages -- a small synthetic near-Earth offset, a larger synthetic offset
combining semi-major-axis and inclination change, and real GTOC4 asteroids -- is used to try to
scale training up from an easy instance towards the full problem, following the same instrument
used throughout: whenever an assumption in the original project plan turned out not to match
reality (a training run failing, a target being physically unreachable in the assumed time
window, background training being interrupted), the discrepancy was diagnosed empirically and
documented rather than silently patched over. Several of the results in Chapter~\ref{chapter:results}
are exactly such documented, negative findings.

\section{Report structure}

Chapter~\ref{chapter:methodology} describes the dynamics model and its validation against Tudat,
the Gymnasium environment, the reward shaping, the curriculum, and the PPO training setup.
Chapter~\ref{chapter:results} reports training outcomes across the curriculum, a sensitivity
analysis of the environment's reward/dynamics knobs, an analysis of the learned stage-1
controller's structure, and an independent verification of that controller's thrust profile
under Tudat. Chapter~\ref{chapter:discussion} discusses what these results say about the two
scientific questions above, and concludes.
